\documentclass[11pt]{article}

\usepackage[margin=1.1in]{geometry}
\usepackage{amsmath,amssymb,amsthm}
\usepackage[numbers]{natbib}
\usepackage{booktabs}
\usepackage[colorlinks=true,linkcolor=blue,citecolor=blue]{hyperref}
\usepackage{microtype}

\newtheorem{theorem}{Theorem}[section]
\newtheorem{proposition}[theorem]{Proposition}
\newtheorem{lemma}[theorem]{Lemma}
\newtheorem{corollary}[theorem]{Corollary}
\theoremstyle{definition}
\newtheorem{definition}[theorem]{Definition}
\newtheorem{assumption}[theorem]{Assumption}
\theoremstyle{remark}
\newtheorem{remark}[theorem]{Remark}

\newcommand{\E}{\mathbb{E}}
\newcommand{\R}{\mathbb{R}}
\newcommand{\Var}{\operatorname{Var}}
\newcommand{\Cov}{\operatorname{Cov}}
\newcommand{\Wt}{\widetilde{W}}
\newcommand{\cut}[1]{c(i) \neq c(j)}

\title{The Retention Chain, in Full:\\
    \large From the Map Equation to Regional Value Loss ---
    Proof Notes for Discussion}
\author{}
\date{July 2026 --- working note, not for distribution}

\begin{document}
\maketitle

% ============================================================
\section{Why this note, and what it claims}
\label{sec:why}

\paragraph{The problem.}
In multi-agent reinforcement learning for traffic signal control (TSC), the
coordination graph is a \emph{design choice}: before training, the $N$
signalized intersections are partitioned into modules, and coordination
machinery (shared rewards, a regional critic, restricted attention) operates
within modules only.
The design question is: \emph{what should a partition preserve, and why?}
Our answer is a three-link chain of inequalities.
Writing $\Phi$ for the fraction of vehicle flow a partition retains within
modules, $\Psi$ for the fraction of inter-agent statistical dependency it
retains, and $J^\star - J_{\mathcal{P}}$ for the value lost by decomposing
the MARL problem regionally:
\begin{center}
    \small
    \begin{tabular}{lll}
        \toprule
        Link                                                        & Statement & Epistemic status \\
        \midrule
        1. map equation $\Leftrightarrow$ flow retention            &
        $q_{\circlearrowleft} = 1 - \Phi$                           &
        exact identity (\S\ref{sec:link1})                                                         \\
        2. cut flow bounds cut dependency                           &
        $1 - \Psi \leq \kappa (1 - \Phi) + \epsilon_{\mathrm{dem}}$ &
        theorem in a stylized model (\S\ref{sec:link2})                                            \\
        3. cut dependency bounds value loss                         &
        $J^\star - J_{\mathcal{P}} \leq C (1 - \Psi) + \epsilon$    &
        imported, informal (\S\ref{sec:link3})                                                     \\
        \bottomrule
    \end{tabular}
\end{center}
Chained, the links say: \emph{a partition that minimizes the map equation
    minimizes an upper bound on the value loss of regional decomposition}.
This is the theoretical backbone of the manuscript; it is why measuring flow
retention --- which requires no training --- is a principled surrogate for
coordination quality.

\paragraph{What this note is for.}
Three things.
(i) Full derivations: every equation is derived from its starting point, so
the argument can be checked line by line.
(ii) Provenance: every imported fact is cited at the point of use.
(iii) An honest strength assessment (\S\ref{sec:strength}), including the
exact places where a critical reader should push.
The note is deliberately self-contained and terse; it is a discussion
document, not a paper section.

\paragraph{One framing remark before starting.}
The chain's central physical object is \emph{influence} --- how much one
intersection's state or action perturbs another's future --- not mutual
information.
Link 2 is proved for an MI instrument because MI is what we can estimate
model-free from data; but MI relates to influence one-sidedly (large observed
MI certifies large influence; small observed MI certifies nothing, since a
strong channel can sit unexercised by the demand distribution).
Keeping this asymmetry in view resolves several apparent puzzles about what
the empirical results do and do not establish; we return to it in
\S\ref{sec:link3}.

% ============================================================
\section{Objects and notation}
\label{sec:notation}

$\mathcal{V}$ is the set of $N$ signalized intersections.
$W_{ij} \geq 0$ is the observed vehicle count traveling $i \to j$ over a
measurement window (from trip logs under a random warm-up policy);
$\Wt_{ij} = W_{ij} + W_{ji}$ is its symmetrization.
A partition $\mathcal{P} = \{P_1, \dots, P_K\}$ assigns each node a module
label $c(i)$.

\begin{definition}[Retention ratios]
    For a non-negative matrix $M$ (flow or dependency), the retention ratio of
    partition $\mathcal{P}$ is
    \begin{equation}
        \Phi(\mathcal{P}, M)
        = \frac{\sum_{k} \sum_{i,j \in P_k} M_{ij}}{\sum_{i,j} M_{ij}},
        \label{eq:retention}
    \end{equation}
    the fraction of total mass whose endpoints share a module.
    We write $\Phi$ when $M$ is flow and $\Psi$ when $M$ is a dependency matrix
    $D$; both lie in $[0,1]$, and $1 - \Phi$ (resp.\ $1 - \Psi$) is the
    \emph{cut} fraction.
\end{definition}

The dependency instrument used in Link 2 is lagged mutual information between
per-intersection states,
$D^{\mathrm{MI}}_{ij} = I\!\left(x_i(t); x_j(t{+}\tau)\right)$, $\tau \geq 1$,
estimated in practice by the KSG estimator~\citep{kraskov2004ksg}.

% ============================================================
\section{Link 1: the map equation is a statement about retained flow}
\label{sec:link1}

\paragraph{Where we start.}
The map equation~\citep{rosvall2008maps,rosvall2009mapequation} scores a
partition $\mathcal{M}$ by the description length of a random walk on the
weighted directed graph:
\begin{equation}
    L(\mathcal{M})
    = q_{\circlearrowleft} H(\mathcal{Q})
    + \sum_{k=1}^{K} p_k^{\circlearrowleft} H(\mathcal{P}_k),
    \label{eq:mapeq}
\end{equation}
where $q_{\circlearrowleft}$ is the total stationary rate of inter-module
steps, $H(\mathcal{Q})$ the entropy of the between-module codebook, and the
second sum prices within-module description.
Infomap searches for the partition minimizing $L$.
The observation --- elementary once seen, but load-bearing --- is that
$q_{\circlearrowleft}$ \emph{is} the cut flow fraction, exactly, under the
walk induced by the empirical flows.

\begin{proposition}[$q_{\circlearrowleft} = 1 - \Phi$]
    \label{prop:phi}
    Define the empirical OD walk by visit rates
    $p_i = \sum_j W_{ij} / S$ with $S = \sum_{a,b} W_{ab}$, and transitions
    $p_{ij} = W_{ij} / \sum_{j'} W_{ij'}$.
    Then
    \begin{equation}
        q_{\circlearrowleft}
        = \sum_{i,j:\, c(i) \neq c(j)} p_i\, p_{ij}
        = 1 - \Phi(\mathcal{P}, W).
    \end{equation}
\end{proposition}

\begin{proof}
    The stationary probability of traversing the directed edge $(i,j)$ in one
    step is
    \[
        p_i\, p_{ij}
        = \frac{\sum_b W_{ib}}{S} \cdot \frac{W_{ij}}{\sum_b W_{ib}}
        = \frac{W_{ij}}{S}.
    \]
    Summing over ordered pairs with $c(i) \neq c(j)$ gives
    \[
        q_{\circlearrowleft}
        = \frac{1}{S} \sum_{c(i) \neq c(j)} W_{ij},
    \]
    which is the complement of Eq.~\eqref{eq:retention}.
\end{proof}

\begin{corollary}[Infomap maximises entropy-regularised flow retention]
    At any fixed codebook-entropy profile, $L(\mathcal{M})$ is increasing in
    $1 - \Phi$; the entropy terms are the regulariser that excludes the
    degenerate one-module solution ($\Phi = 1$ at full within-module
    description cost).
\end{corollary}

\paragraph{Provenance and caveats.}
The identity is close in spirit to the flow-based reading of community
detection formalized by Markov
stability~\citep{delvenne2010stability,lambiotte2014multiscale}: up to a null
model, retaining flow is retaining the one-step clustered autocovariance of
the walk.
Three gaps separate the identity from Infomap-as-software, and all three are
disclosed in the manuscript: real trips are path-dependent while the walk is
memoryless (higher-order map equations are the
refinement~\citep{rosvall2014memory}); directed Infomap smooths visit rates
with teleportation, so its internal walk differs slightly from the empirical
OD walk above; and Infomap's search is greedy over an NP-hard
landscape~\citep{edler2023community}.
None of these threaten the identity itself; they mean the empirical margin of
Infomap over alternatives is an object worth measuring rather than a
foregone conclusion.

% ============================================================
\section{Link 2: cut flow bounds cut dependency (proved in a stylized model)}
\label{sec:link2}

\subsection{Why a stylized model, and which one}
\label{sec:whymodel}

The general claim --- flow cut by a partition upper-bounds the statistical
dependency it cuts --- requires relating an information-theoretic quantity
(MI) to a physical one (flow) through the dynamics.
For general nonlinear stochastic dynamics this is the territory of strong
data-processing inequalities on Bayesian networks, where end-to-end MI is
bounded by sums over paths of products of per-edge contraction
coefficients~\citep{evans1999percolation,polyanskiy2017sdpi}; making the
per-edge coefficients explicit functions of traffic flow is open.
We therefore prove the bound in the model where every step is exact: the
\emph{linearization of a link-queue / cell-transmission
    model}~\citep{daganzo1994ctm,jin2012linkqueue} about an undersaturated
operating point.
This is the standard stylisation in traffic flow theory, and --- the point of
choosing it --- the manuscript's three assumptions hold in it \emph{literally}
rather than approximately.

Let $x(t) \in \R^N$ collect per-intersection queue deviations from the
operating point, and posit
\begin{equation}
    x(t+1) = A\, x(t) + b\, \xi(t) + w(t),
    \label{eq:var}
\end{equation}
with the following conditions and their traffic readings:

\begin{itemize}
    \item[(M1)] \textbf{Locality.} $A_{ij} \neq 0$ only if
          $j \in \mathcal{N}(i) \cup \{i\}$.
          \emph{Reading:} one-step queue evolution depends on neighboring
          queues only --- the support structure of CTM/link-queue
          dynamics~\citep{daganzo1994ctm,jin2012linkqueue}.
    \item[(M2)] \textbf{Flow-bounded coupling.} $|A_{ij}| \leq \alpha
              \Wt_{ij}$ for $i \neq j$.
          \emph{Reading:} the linearised sensitivity of queue $i$ to queue $j$
          is carried by the two physical channels --- forward (arrivals routed
          $j \to i$) and backward (blocking / spill-back, which propagates
          against the flow direction, hence the symmetrization $\Wt$) --- each
          proportional to the flow on the link, with $\alpha$ absorbing
          capacities and saturation rates.
          Common-demand coupling is carried by $b$, \emph{not} by $A$.
    \item[(M3)] \textbf{Sub-criticality.}
          $\max\!\left(\| |A| \|_1, \| |A| \|_\infty\right) \leq 1 - \delta$
          for some $\delta \in (0,1]$, where $|A|$ is the entrywise absolute
          value.
          \emph{Reading:} the total coupling any intersection exerts on its
          neighbors, and receives from them, is strictly below unity --- the
          undersaturated regime, in which queues are stable.
          (Both norms are needed below; requiring both is the honest price of
          the walk-counting argument.)
    \item[(M4)] \textbf{Noise and demand.} $w(t) \overset{\text{iid}}{\sim}
              \mathcal{N}(0, \Sigma_w)$ with $\Sigma_w$ diagonal,
          $\sigma_{\min}^2 \leq (\Sigma_w)_{ii} \leq \sigma_{\max}^2$;
          $\xi(t) \overset{\text{iid}}{\sim} \mathcal{N}(0, \sigma_\xi^2)$
          independent of $w$.
          \emph{Reading:} idiosyncratic per-intersection fluctuations with a
          noise floor, plus one scalar network-wide demand factor --- the
          common cause that will become $\epsilon_{\mathrm{dem}}$.
\end{itemize}

Time-invariance of $(A, b, \Sigma_w, \sigma_\xi)$ is the model's version of
regime stationarity.
Note $\rho(A) \leq \rho(|A|) \leq \| |A| \|_\infty \leq 1 - \delta < 1$, so a
stationary causal solution exists.

\subsection{Covariance structure: where the matrices come from}
\label{sec:cov}

\paragraph{Lagged covariance.}
Iterating Eq.~\eqref{eq:var} $\tau$ steps,
\begin{equation}
    x(t + \tau)
    = A^\tau x(t)
    + \sum_{k=0}^{\tau - 1} A^{k}\!\left(b\, \xi(t + \tau - 1 - k)
    + w(t + \tau - 1 - k)\right).
    \label{eq:iterate}
\end{equation}
Every random term in the sum is dated $\geq t$ and is therefore independent
of $x(t)$ (which, in the stationary causal solution, is a function of
randomness dated $< t$).
Hence
\begin{equation}
    \Gamma_\tau := \Cov\!\left(x(t{+}\tau),\, x(t)\right) = A^\tau \Sigma,
    \qquad \Sigma := \Var\!\left(x(t)\right).
    \label{eq:gamma}
\end{equation}

\paragraph{Stationary covariance and its two parts.}
Taking variances in Eq.~\eqref{eq:var} and using independence of
$x(t), \xi(t), w(t)$:
$\Sigma = A \Sigma A^{\!\top} + \sigma_\xi^2\, b b^{\!\top} + \Sigma_w$,
the discrete Lyapunov equation, whose unique solution under
$\rho(A) < 1$ is the convergent series~\citep[Ch.~5]{horn2013matrix}
\begin{equation}
    \Sigma
    = \underbrace{\sum_{k \geq 0} A^{k} \Sigma_w (A^{\!\top})^{k}}_{=:\,
        \Sigma_0 \ \text{(noise-driven)}}
    \;+\;
    \underbrace{\sigma_\xi^2 \sum_{k \geq 0} A^{k} b\, b^{\!\top}
            (A^{\!\top})^{k}}_{=:\, \Sigma_\xi \ \text{(demand-driven)}}.
    \label{eq:lyap}
\end{equation}
The split is the model's causal anatomy: $\Sigma_0$ is dependency generated
by idiosyncratic noise propagating through the network (flow-borne, the
channel we want to bound by flow), $\Sigma_\xi$ is dependency generated by
the common demand factor (the confounder, which no flow bound can and should
control).
Two facts for later: $\Sigma \succeq \Sigma_w \succeq \sigma_{\min}^2 I$, and
$\bar{\sigma}^2 := \max_i \Sigma_{ii} < \infty$.

\subsection{Cut charging: every cross-module correlation pays a cut-edge toll}
\label{sec:cutcharge}

\paragraph{Motivation.}
The dependency between $x_i(t)$ and $x_j(t{+}\tau)$ in this model is
generated by noise injected at some intersection $m$ propagating to both $i$
and $j$.
(Note this includes, and correctly prices, the ``third-party'' channel $m
    \notin \{i, j\}$ that informal path arguments miss.)
If $i$ and $j$ sit in different modules, at least one of the two propagation
walks must cross the module boundary --- and (M2) prices every crossing step
by flow.
That is the entire idea; the lemma is bookkeeping.

\begin{lemma}[Cut charging]
    \label{lem:cutcharge}
    For every partition $\mathcal{P}$ and every $\tau \geq 1$,
    \begin{equation}
        \sum_{c(i) \neq c(j)}
        \left|\left(A^\tau \Sigma_0\right)_{ji}\right|
        \;\leq\;
        \frac{2\, \sigma_{\max}^2\, \alpha}{\delta^{3}}
        \sum_{c(u) \neq c(v)} \Wt_{uv}.
        \label{eq:cutcharge}
    \end{equation}
\end{lemma}

\begin{proof}
    \emph{Step 1: walk expansion.}
    From Eq.~\eqref{eq:lyap}, using diagonality of $\Sigma_w$,
    \begin{equation}
        (A^\tau \Sigma_0)_{ji}
        = \sum_{k \geq 0} \sum_{m}
        (A^{\tau + k})_{jm}\, (\Sigma_w)_{mm}\, (A^{k})_{im}.
        \label{eq:walkpair}
    \end{equation}
    Bounding entrywise by $|A|$, each term pairs a walk $m \to i$ of length
    $k$ with a walk $m \to j$ of length $\tau + k$ in the coupling graph
    (entry $(p, q)$ of $|A|^{\ell}$ sums the weight products
    $\prod_e |A_e|$ over all length-$\ell$ walks $q \to p$).

    \emph{Step 2: one walk must cross.}
    Fix a cut pair, $c(i) \neq c(j)$.
    If $c(m) = c(i)$, the walk $m \to j$ has endpoints in different modules
    and hence uses at least one crossing step; its length $\tau + k \geq 1$
    since $\tau \geq 1$.
    If $c(m) \neq c(i)$, the walk $m \to i$ crosses; its length $k$ cannot be
    $0$, because $k = 0$ forces $m = i$ (the $(|A|^0)_{im} = \delta_{im}$
        term), contradicting $c(m) \neq c(i)$.
        So in all cases at least one of the two walks crosses.

        \emph{Step 3: first-crossing decomposition.}
        Decompose any crossing walk at its \emph{first} crossing step
    $(u \to v)$, $c(u) \neq c(v)$: the walk splits uniquely into a prefix
        (start node to $u$, arbitrary length $\ell_1 \geq 0$), the crossing edge,
        and a suffix ($v$ to end node, arbitrary length $\ell_2 \geq 0$).
        The crossing edge's weight obeys $|A_{vu}| \leq \alpha \Wt_{uv}$ by (M2).
        Dropping the constraints that the prefix not cross and that
    $\ell_1 + 1 + \ell_2$ equal the original length only enlarges the sum.

        \emph{Step 4: resolvent sums.}
        All quantities are non-negative, so coupled sums decouple,
        \[
        \textstyle
        \sum_k a_k\, b_{\tau + k}
        \;\leq\;
        \big(\sum_k a_k\big)\big(\sum_\ell b_\ell\big).
    \]
    By the Neumann series~\citep[Ch.~5]{horn2013matrix},
    $\sum_{\ell \geq 0} |A|^{\ell} = (I - |A|)^{-1}$, whose row sums are at
    most $1/\delta$ (from $\| |A| \|_\infty \leq 1 - \delta$) and whose column
    sums are at most $1/\delta$ (from $\| |A| \|_1 \leq 1 - \delta$).
    Now take the case in which the crossing walk is $m \to j$ (the other case
    is symmetric).
    Summing the non-crossing walk over its free endpoint $i$ and all lengths
    $k$ gives a column sum $\leq 1/\delta$, uniformly in $m$.
    Summing the crossing walk over $m$, $j$, and all lengths gives, via
    Step 3,
    \[
        \sum_{(u,v):\, c(u) \neq c(v)} \alpha \Wt_{uv}
        \underbrace{\Big(\sum_{m, \ell_1} (|A|^{\ell_1})_{um}\Big)}_{\text{row
                sum} \, \leq\, 1/\delta}
        \underbrace{\Big(\sum_{j, \ell_2} (|A|^{\ell_2})_{jv}\Big)}_{\text{column
                sum} \, \leq\, 1/\delta}.
    \]
    Collecting: each case contributes at most
    $\sigma_{\max}^2\, \alpha\, \delta^{-3} \sum_{\mathrm{cut}} \Wt_{uv}$, and
    there are two cases.
\end{proof}

\subsection{From correlation to mutual information}
\label{sec:mi}

\paragraph{Motivation.}
Lemma~\ref{lem:cutcharge} bounds \emph{covariances}; the instrument is
\emph{MI}.
For jointly Gaussian pairs the two are in bijection,
$I = -\tfrac12 \ln(1 - \rho^2)$~\citep[Ch.~8]{cover2006elements}, but the
bijection is nonlinear and blows up as $|\rho| \to 1$.
The noise floor in (M4) is exactly what caps $|\rho|$ away from $1$ and makes
MI \emph{linear} in correlation up to a constant.

\begin{lemma}[MI is linear in correlation under a noise floor]
    \label{lem:gaussmi}
    Let $\bar{\rho}^2 := 1 - \sigma_{\min}^2 / \bar{\sigma}^2$.
    For all pairs $(i,j)$ and lags $\tau \geq 1$:
    $|\rho_{ij}(\tau)| \leq \bar{\rho} < 1$, and
    \begin{equation}
        D^{\mathrm{MI}}_{ij}
        = -\tfrac{1}{2}\ln\!\big(1 - \rho_{ij}^2(\tau)\big)
        \;\leq\;
        \frac{\bar{\rho}}{2(1 - \bar{\rho}^2)}\, |\rho_{ij}(\tau)|
        \;\leq\;
        \frac{\bar{\rho}}{2(1 - \bar{\rho}^2)\, \sigma_{\min}^2}
        \big|(\Gamma_\tau)_{ji}\big|.
        \label{eq:milinear}
    \end{equation}
\end{lemma}

\begin{proof}
    \emph{Correlation cap.}
    For $\tau \geq 1$, $x_i(t)$ is measurable with respect to the history up
    to time $t + \tau - 1$.
    For jointly Gaussian variables, conditioning on more variables cannot
    increase conditional variance, so
    \[
        \Var\!\left(x_j(t{+}\tau) \mid x_i(t)\right)
        \;\geq\;
        \Var\!\left(x_j(t{+}\tau) \mid \text{history to } t{+}\tau{-}1,\,
        \xi(t{+}\tau{-}1)\right)
        = (\Sigma_w)_{jj} \geq \sigma_{\min}^2,
    \]
    the innovation variance, by Eq.~\eqref{eq:var}.
    Since for a bivariate Gaussian
    $1 - \rho^2 = \Var(Y \mid X) / \Var(Y)$ and
    $\Var(x_j) \leq \bar{\sigma}^2$, we get
    $1 - \rho_{ij}^2 \geq \sigma_{\min}^2 / \bar{\sigma}^2$, i.e.\
    $|\rho_{ij}| \leq \bar{\rho}$.
    (This is where $\tau \geq 1$ is used; at $\tau = 0$ the innovation
    argument fails.)

    \emph{Secant bound.}
    $u \mapsto -\ln(1 - u)$ is convex and vanishes at $0$, so on
    $[0, \bar{u}]$ it lies below its chord:
    $-\ln(1 - u) \leq \frac{u}{\bar{u}}\big(-\ln(1 - \bar{u})\big)
        \leq \frac{u}{1 - \bar{u}}$,
    the last step because
    $-\ln(1 - \bar{u}) = \ln\!\big(1 + \tfrac{\bar{u}}{1 - \bar{u}}\big) \leq
        \tfrac{\bar{u}}{1 - \bar{u}}$.
    Apply with $u = \rho^2 \leq \bar{u} = \bar{\rho}^2$ and use
    $\rho^2 \leq \bar{\rho}\,|\rho|$ to obtain the first inequality of
    Eq.~\eqref{eq:milinear}.
    The second follows from
    $|\rho_{ij}| = |(\Gamma_\tau)_{ji}| / (\mathrm{sd}_i\, \mathrm{sd}_j)$ and
    $\mathrm{sd}_i, \mathrm{sd}_j \geq \sigma_{\min}$
    (since $\Sigma \succeq \sigma_{\min}^2 I$).
\end{proof}

\subsection{The transfer theorem}
\label{sec:transferthm}

\begin{theorem}[Cut flow bounds cut dependency; stylized model]
    \label{thm:transfer}
    Under (M1)--(M4), for every partition $\mathcal{P}$ and every lag
    $\tau \geq 1$,
    \begin{equation}
        1 - \Psi(\mathcal{P}, D^{\mathrm{MI}})
        \;\leq\;
        \kappa \left(1 - \Phi(\mathcal{P}, \Wt)\right)
        + \epsilon_{\mathrm{dem}},
        \label{eq:transfer}
    \end{equation}
    with explicit constants
    \begin{equation*}
        \kappa
        = \frac{\bar{\rho}\, \sigma_{\max}^2\, \alpha}
        {(1 - \bar{\rho}^2)\, \sigma_{\min}^2\, \delta^{3}}
        \cdot
        \frac{\sum_{u \neq v} \Wt_{uv}}{\sum_{i \neq j} D^{\mathrm{MI}}_{ij}},
        \qquad
        \epsilon_{\mathrm{dem}}
        = \frac{\bar{\rho}}{2(1 - \bar{\rho}^2)\, \sigma_{\min}^2}
        \cdot
        \frac{\sum_{c(i) \neq c(j)}
            \big|(A^\tau \Sigma_\xi)_{ji}\big|}{\sum_{i \neq j}
            D^{\mathrm{MI}}_{ij}}.
    \end{equation*}
    The lag-averaged instrument
    $\bar{D}^{\mathrm{MI}} = \tfrac{1}{|T|}\sum_{\tau \in T} D^{\mathrm{MI}}(\tau)$
    obeys the same bound with the total mass in the constants replaced by that
    of $\bar{D}^{\mathrm{MI}}$, since Eq.~\eqref{eq:transfer}'s unnormalized
    form holds uniformly over $\tau \geq 1$.
\end{theorem}

\begin{proof}
    Split $\Gamma_\tau = A^\tau \Sigma_0 + A^\tau \Sigma_\xi$
    (Eqs.~\eqref{eq:gamma}--\eqref{eq:lyap}).
    Apply Lemma~\ref{lem:gaussmi} to each cut pair and sum:
    \[
        \sum_{c(i) \neq c(j)} D^{\mathrm{MI}}_{ij}
        \;\leq\;
        \frac{\bar{\rho}}{2(1 - \bar{\rho}^2)\sigma_{\min}^2}
        \left(
        \sum_{c(i) \neq c(j)} \big|(A^\tau \Sigma_0)_{ji}\big|
        + \sum_{c(i) \neq c(j)} \big|(A^\tau \Sigma_\xi)_{ji}\big|
        \right).
    \]
    Bound the first bracket by Lemma~\ref{lem:cutcharge}; leave the second as
    the numerator of $\epsilon_{\mathrm{dem}}$.
    Divide by the total dependency mass $\sum_{i \neq j} D^{\mathrm{MI}}_{ij}$
    and use
    $1 - \Phi = \sum_{\mathrm{cut}} \Wt_{uv} / \sum_{u \neq v} \Wt_{uv}$;
    the factor $2$ from the lemma combines with the $2$ in
    Lemma~\ref{lem:gaussmi}'s constant to give the stated $\kappa$.
\end{proof}

\begin{corollary}[Demand conditioning removes the confounder exactly]
    \label{cor:demzero}
    Conditioned on the demand path $\{\xi_s\}$, the process~\eqref{eq:var} has
    the same law with $\sigma_\xi = 0$, so $\Sigma_\xi = 0$ and the
    conditional-MI instrument satisfies Eq.~\eqref{eq:transfer} with
    $\epsilon_{\mathrm{dem}} = 0$.
    Operationally, aggregate network occupancy proxies $\xi$.
\end{corollary}

\paragraph{Interpretation.}
The direction of Theorem~\ref{thm:transfer} matters: low cut flow
\emph{certifies} low cut dependency (up to the demand term), but high
retention does not certify high task performance --- the bound is a necessary
condition for lossless decomposition, not a sufficient one.
The constant's $\delta^{-3}$ makes a second point quantitative: as the
network approaches saturation ($\delta \to 0$), the bound degrades --- and
saturation is precisely where the backward (spill-back) channel activates.
The theory is thus most informative in the undersaturated regime and honest
about weakening where coordination pressure is greatest.

\subsection{How strong is Link 2?}
\label{sec:link2strength}

\emph{Rigorous:} everything above, within the model.
No step is heuristic; the constants are explicit; the confounder is isolated
structurally (in $b$) rather than assumed small; and
Corollary~\ref{cor:demzero} converts the manuscript's ``conditional-MI fix''
from a proposal into a theorem.

\emph{Idealisations a critic should probe} (each is a genuine limitation; we
state our best defense):
\begin{enumerate}
    \item \textbf{Linearity.} Eq.~\eqref{eq:var} is a linearization; real
          queue dynamics are piecewise (free-flow vs.\ congested branches of
          the fundamental diagram).
          Defense: the linearization is standard for undersaturated analysis,
          and the general nonlinear route is identified --- SDPI/information
          percolation~\citep{evans1999percolation,polyanskiy2017sdpi} gives
          the same walk-sum structure with per-edge contraction coefficients
          replacing $|A_e|$; what is missing is only the traffic-physics step
          bounding those coefficients by flow.
    \item \textbf{Gaussianity.} Used twice: the MI--correlation bijection and
          conditional-variance monotonicity.
          Defense: both have sub-Gaussian analogues at the price of constants;
          nothing in the walk-counting core uses Gaussianity.
    \item \textbf{Dual-norm sub-criticality (M3).} Row sums bound influence
          \emph{exerted}; column sums bound influence \emph{received}.
          Both below $1 - \delta$ is stronger than mere stability
          ($\rho(A) < 1$).
          Defense: physically, an intersection neither exerts nor receives
          unbounded total coupling in an undersaturated network; but this is
          the cleanest place to attack the constants.
    \item \textbf{Scalar states, i.i.d.\ demand, $\tau \geq 1$.} Scalar
          per-agent states match the queue instruments (vector states extend
          via the log-determinant Gaussian MI at notational cost); an AR
          demand factor changes $\Sigma_\xi$ but not the split; $\tau = 0$ is
          genuinely excluded (the innovation argument fails), and the
          empirical instrument uses $\tau \in \{1, 5, 15, 30\}$ only.
    \item \textbf{Tightness.} The bound is a certificate, not an equality;
          the $\delta^{-3}$ is likely loose (each of the three resolvent sums
          is bounded crudely).
          A numerical companion --- simulate Eq.~\eqref{eq:var} on the real
          Cologne8/Ingolstadt21 graphs with measured $W$, compute exact
          Gaussian MI from Eq.~\eqref{eq:lyap}, overlay the bound across all
          feasible benchmark partitions --- would measure the slack directly
          and is planned but not yet run.
\end{enumerate}

% ============================================================
\section{Link 3: cut dependency bounds regional value loss (imported)}
\label{sec:link3}

\paragraph{Where the machinery comes from.}
This link is not proved in this note or in the manuscript; it is imported
from two literatures, and the import is the link's weakness.
(i) \emph{Networked MARL with decaying influence}: when one-step
inter-agent influence is local, the sensitivity of agent $i$'s value to agent
$j$ decays exponentially in graph distance, and policies truncated to local
neighborhoods approximate the global optimum with bias controlled by the
influence discarded~\citep{qu2020average,qu2022scalable}; the classical
row-sum (Dobrushin-type) condition for this can be weakened to a
policy-dependent spectral certificate~\citep{chakraborty2026locality}.
(ii) \emph{Influence-based abstraction}: the value loss of a local model is
bounded by the error of its representation of the influence the rest of the
system exerts on it~\citep{congeduti2021iba}.
A regional decomposition is such a truncation --- each module models
within-module influence and discards the rest --- so the discarded quantity
is exactly the cut influence mass, giving the shape
\begin{equation}
    J^\star - J_{\mathcal{P}}
    \;\leq\;
    C \left(1 - \Psi(\mathcal{P}, D)\right)
    + \epsilon_{\mathrm{approx}} + \epsilon_{\mathrm{opt}},
    \label{eq:valueloss}
\end{equation}
with $D$ an \emph{influence} matrix and the $\epsilon$'s collecting
function-approximation and optimization error the partition does not control.

\paragraph{Two formal gaps, stated precisely.}
First, the cited neighborhood-truncation guarantees are proved for specific
algorithms and settings ($\kappa$-hop truncated policies/critics,
average-reward or discounted variants); a module is a truncation to a
\emph{set}, not a $\kappa$-hop ball, so the set-based IBA
route~\citep{congeduti2021iba} is the closer parent, and writing out
Eq.~\eqref{eq:valueloss} as a theorem for regional MAPPO under (M1)--(M4) is
open work.
(In the stylized linear model with quadratic queue cost this becomes a
decentralized-LQR question with block-structured information --- plausibly
tractable, and in our view the highest-value theory extension.)
Second, Eq.~\eqref{eq:valueloss} wants \emph{influence}; the measurable
instrument is \emph{MI}, and the two relate one-sidedly: by data processing,
MI observable across a cut is upper-bounded by (a function of) the influence
crossing it, so large measured cut-MI certifies large cut influence, while
small measured cut-MI certifies nothing.
$\Psi(D^{\mathrm{MI}})$ comparisons therefore \emph{rule partitions out} (a
partition cutting much MI provably cuts much influence, so its bound
\eqref{eq:valueloss} is poor) rather than certifying winners.
This one-sidedness is not a flaw to hide; it is the correct epistemic status
of every dependency-retention comparison in the manuscript, and stating it
preempts the sharpest referee question (``what does high $\Psi$ actually
establish?'': by itself, nothing --- it is the \emph{low}-$\Psi$ partitions
about which the theory speaks).

\paragraph{Empirical status of the chain's causal content.}
The state-level instrument ($D^{\mathrm{MI}}$) shows the predicted retention
orderings across seeds.
The action-level (causal) content has now been probed twice: observationally
(conditional MI of actions on future queues; no edge survives FDR in 60
tests) and --- decisively for identification --- by a preregistered
\emph{randomized action intervention} on the frozen policies, in which the
logged random assignment is exogenous by construction.
The randomized design passed its calibration checks (negative controls at
nominal false-positive rates; manipulation check on action entropy) and
\emph{still} produced zero FDR-significant held-out edges on either network,
with effect estimates at sub-vehicle scale, though only $14/22$ and $41/71$
preregistered pairs reached prospective $80\%$ power.
The causal reading of Link 3 is therefore tested-and-null under the current
demand regime, not merely unresolved.
Three readings survive: closed-loop damping (perturbations measured under
trained stabilising policies at all other intersections, which attenuate
disturbances by design); a deeply undersaturated regime (converging with the
empty spill-back diagnostic, and consistent with the transfer bound's
$\delta^{-3}$ locating strong coupling near saturation); and residual power.
The constructive consequence: if cuttable influence is near zero for every
partition, Eq.~\eqref{eq:valueloss}'s coupling term cannot distinguish
partitions, and the observed large partition-dependent training differences
fall to $\epsilon_{\mathrm{opt}}$ --- \emph{optimization conditioning} ---
for which correlated within-module signals (the confirmed state-MI retention,
common-demand mass included) suffice; causation is not required for the
variance-reduction mechanisms of reward sharing and regional critics.

% ============================================================
\section{Epistemic summary, and where to push}
\label{sec:strength}

\begin{center}
    \small
    \begin{tabular}{p{2.6cm} p{3.4cm} p{4.2cm} p{3.6cm}}
        \toprule
        Link & Status                                           & Main attack surface
             & Remedy / response                                                      \\
        \midrule
        1 (map eq.\ $=$ flow)
             & Exact identity                                   &
        Teleportation; memoryless walk; greedy search
             & Disclosed; empirical margins quantify the slack                        \\
        \addlinespace
        2 (flow $\Rightarrow$ dependency)
             & Theorem, stylized model                          &
        Linearity; Gaussianity; dual-norm (M3); $\delta^{-3}$ looseness
             & SDPI route for general case; numerical companion                       \\
        \addlinespace
        3 (dependency $\Rightarrow$ value loss)
             & Imported, informal                               &
        Algorithm-specificity; module $\neq$ $\kappa$-hop ball; MI vs.\
        influence one-sidedness
             & IBA set-based route; decentralized-LQR version;
        interventional CMI                                                            \\
        \bottomrule
    \end{tabular}
\end{center}

\paragraph{Checklist for a critical reading} (the six places we would push,
in descending order of consequence):
\begin{enumerate}
    \item Link 3's import: does the set-based IBA bound actually deliver
          Eq.~\eqref{eq:valueloss} for a regional critic, or only for a
          regional \emph{model}? (\S\ref{sec:link3})
    \item (M3)'s column-sum half: is bounded total \emph{received} coupling
          defensible near merges/diverges in real networks?
          (\S\ref{sec:whymodel})
    \item The first-crossing decomposition and the decoupling inequality
          $\sum_k a_k b_{\tau+k} \leq (\sum a)(\sum b)$: both are correct but
          crude; is there a tighter charging scheme? (\S\ref{sec:cutcharge})
    \item The conditional-variance cap and its $\tau \geq 1$ restriction:
          any instrument at $\tau = 0$ is outside the theorem.
          (\S\ref{sec:mi})
    \item $\epsilon_{\mathrm{dem}}$'s growth with network size: the model
          predicts, and the data confirm, weaker raw flow--MI alignment on
          the larger network; Corollary~\ref{cor:demzero} is the fix but has
          not yet been exercised empirically. (\S\ref{sec:transferthm})
    \item The Gaussian MI formula's role: KSG estimates on real traces are
          not Gaussian MI; the theorem speaks about the model's MI, and the
          transfer of the \emph{ordering} prediction to KSG estimates is an
          empirical, not logical, step. (\S\ref{sec:mi})
\end{enumerate}

\paragraph{Bottom line.}
Link 1 is unimpeachable but near-known; its value is the connection it
anchors.
Link 2 is now a real theorem in a model whose assumptions are the
manuscript's assumptions made literal --- the strongest formal component, and
the right object to present for proof-checking.
Link 3 is an interpretation with named parents and two named gaps; it should
be defended as such, never as a proved result.
The chain's overall claim --- map-equation minimization minimizes an upper
bound on regional value loss --- is accordingly: exact at step 1, proved in a
stylized model at step 2, and conditional on imported machinery at step 3.

    % ============================================================
    {\small
        \bibliographystyle{plainnat}
        \bibliography{refs}
    }

\end{document}
